% ID: SEL-D0029-OB-E01
% ESTADO: validada. S030, cotejada con examen y criterios.
Sean $x,y$ las cantidades de palas A y B. Al pasar los kilos
a gramos, las cotas de fibra y goma se simplifican a
$9x+10y\le750$ y $2x+y\le130$. También $0\le x\le60$ e $y\ge0$.
En el eje vertical limita la fibra y se obtiene $(0,75)$.
Sobre $x=60$ limita la goma y se obtiene $(60,10)$.
Las dos fronteras de materiales se cruzan en $(50,30)$.
\begin{center}
\setlength{\unitlength}{7mm}
\begin{picture}(10,10)(-2,-1)
\qbezier(0,0)(3.5,0)(7,0)\put(7.2,0){$x$}
\qbezier(0,0)(0,4)(0,8)\put(0,8.2){$y$}
\qbezier(0,7.5)(2.5,5.25)(5,3)
\qbezier(5,3)(5.5,2)(6,1)
\qbezier(6,1)(6,0.5)(6,0)
\put(0,7.5){\makebox(0,0)[r]{$(0,75)$}}
\put(5.2,3.2){$(50,30)$}
\put(6.2,1.1){$(60,10)$}
\put(6,-0.4){\makebox(0,0)[t]{$(60,0)$}}
\end{picture}
\end{center}
El beneficio lineal $30x+20y$ toma los valores $0,1800,2000,2100,1500$
en $(0,0)$, $(60,0)$, $(60,10)$, $(50,30)$ y $(0,75)$.
El mejor vértice tiene coordenadas enteras. Se deben fabricar
\textbf{50 palas A y 30 palas B}, con $\boxed{2100}$ euros de beneficio.
La propuesta $(49,32)$ requiere $90(49)+100(32)=7610$ g de fibra,
por encima de 7500 g, de modo que es inviable.
